\begin{titlepage}
  \centering
  \vspace*{8mm}
  {\bfseries\fontsize{26}{31}\selectfont BrainChip AKD1000\par}
  \vspace{5mm}
  {\fontsize{16}{21}\selectfont
  Bring-up, model execution and troubleshooting manual\par}
  \vspace{8mm}
  \includegraphics[width=0.60\textwidth]{hardware-powered-isolated.png}
  \vfill
  {\large Álvaro Poveda Ruiz\par}
  \vspace{2mm}
  {\small Supplementary material for the Bachelor's Thesis (TFG)\par
  Universidad Politécnica de Madrid, 2026\par}
  \vspace{8mm}
  {\small Version \versionmanual{} · Revision: \fechacorte\par}
\end{titlepage}

\chapter*{Abstract}
\addcontentsline{toc}{chapter}{Abstract}

This manual describes the bring-up of a BrainChip AKD1000 board connected over PCIe.
The procedure covers host preparation, driver and Python environment installation,
model loading and execution checks. It also explains how to measure latency and
identify common faults.

The first test uses a simple model with a known output. It runs in software first
and then on the board. This checks the Akida environment, model compatibility and
hardware access separately before running your own application.

Appendix~\ref{app:tfg} documents the thesis setup, software versions and results.
The included tools can repeat the smoke test, but do not reproduce the AoA experiment
on their own: its models and data are not part of this release. A complete run of this
version of the procedure on a physical board remains pending.

\section*{Scope}

The procedure covers environment installation, board detection, model loading and
execution, statistics and performance measurement. It also includes checks for
connection, driver, data and model errors.

Apply the assembly instructions to the specific revision of each component.
Before connecting power, check the voltage, polarity and cable orientation in the
manufacturer's documentation. Board repair and internal electronic design are outside
the scope of this manual.

\chapter*{How to use this document}
\addcontentsline{toc}{chapter}{How to use this document}

Download the PDF together with \code{akd1000-bringup.zip} and extract the ZIP.
Open a terminal in \projectroot, which contains \code{pyproject.toml},
\code{examples/} and \code{models/}. All project commands start from this directory
unless another location is specified. Do not run the examples from the ZIP viewer.
Keep the PDF and ZIP at the same version.

Run only the block for your host. Blocks that use
\code{sudo}, \code{apt} or \code{lspci} are for Linux. \Cref{chap:software}
includes environment activation on Windows.

Without a board, start at \cref{chap:software} and complete the software test.
With hardware available, first read the inventory and assembly steps in chapters
\ref{chap:orientacion} and \ref{chap:hardware}, then prepare the driver and SDK.

\section*{Version and included material}

\begin{tabularx}{\textwidth}{@{}L{0.29\textwidth}Y@{}}
\toprule
\textbf{Field} & \textbf{Value} \\
\midrule
Version and date & \versionmanual{}; \fechacorte \\
Editable source & \code{manual/main.tex} and \code{manual/sections/} \\
Thesis reference & Python 3.11.2 and Akida 2.19.1; appendix~\ref{app:tfg} \\
Tests for this edition & Automated tests and simulator with Akida 2.19.1 and 2.19.3 on Windows; physical execution pending \\
Included model & \code{models/akida\_v1\_smoke\_test.fbz}; four inputs and two outputs, no AoA data \\
\bottomrule
\end{tabularx}

The bibliography links to the official instructions. The current software route and
driver compatibility are checked separately: installing an SDK does not guarantee
that the kernel can load the PCIe driver.
